% !TEX root = ../../ctfp-print.tex

\lettrine[lhang=0.17]{幺半群}{是范畴论}和编程中都很重要的概念. 范畴对应强类型语言, 幺半群对应无类型
语言. 这是因为在幺半群里你可以组合任意两个箭头, 就像在无类型语言里你可以组合任意两个函数一样
(当然, 执行程序时你可能最终会得到一个运行时错误).

我们已经看到, 幺半群可以被描述成一个只有一个物件的范畴, 其中所有逻辑都编码在态射组合的规则里.
这个范畴模型与更传统的集合论定义完全等价 --- 在集合论定义里, 我们把两个元素 ``相乘'' 得到第三个
元素. 这个 ``乘法'' 过程还可以进一步拆解: 先形成一个元素的配对, 再把这个配对与某个既有元素
(它们的 ``积'') 等同起来.

如果我们放弃乘法的第二部分 --- 把配对与既有元素等同起来 --- 会发生什么? 我们可以从一个任意的集合
出发, 形成所有可能的元素配对, 并把它们叫作新元素. 然后我们再把这些新元素与所有可能的元素配对, 如此
往复. 这是一场链式反应 --- 我们会永远不断地添加新元素. 最终得到的无穷集合 \emph{几乎} 是个幺半群.
但幺半群还需要一个单位元素和结合律. 没问题, 我们可以添加一个特殊的单位元素, 并等同其中的一些配对 ---
恰好足够满足单位律和结合律即可.

看一个简单的例子. 我们从只有两个元素的集合 $\{a, b\}$ 出发. 我们叫它们自由幺半群的\emph{生成元}.
首先, 我们添加一个特殊元素 $e$ 充当单位. 接下来我们添加所有元素的配对, 叫它们 ``积''.
$a$ 和 $b$ 的积就是配对 $(a, b)$. $b$ 和 $a$ 的积是配对
$(b, a)$, $a$ 与 $a$ 的积是 $(a, a)$, $b$ 与 $b$ 的积是
$(b, b)$. 我们也可以形成与 $e$ 的配对, 比如 $(a, e)$、
$(e, b)$ 等等, 但我们把它们分别与 $a$、$b$ 等同起来. 所以这一轮我们只添加了
$(a, a)$、$(a, b)$、$(b, a)$ 和 $(b, b)$, 最终得到集合
$\{e, a, b, (a, a), (a, b), (b, a), (b, b)\}$.

\begin{figure}[H]
  \centering
  \includegraphics[width=0.8\textwidth]{images/bunnies.jpg}
\end{figure}

\noindent
下一轮我们会继续添加这样的元素: $(a, (a, b))$、$((a, b), a)$ 等等.
此时我们必须保证结合律成立, 所以会把 $(a, (b, a))$ 与 $((a, b), a)$
等同起来, 以此类推. 换句话说, 我们不再需要内部的括号.

你可以猜到这个过程的最终结果: 我们会创造出 $a$ 和 $b$ 的所有可能的列表. 事实上, 如果把
$e$ 表示成空列表, 就能看出我们的 ``乘法'' 不过是列表拼接而已.

这类构造 --- 不断生成元素的所有可能组合, 并且只做最少量的等同 --- 恰好足以维持那些定律 ---
叫作\emph{自由}构造. 我们刚才所做的,就是从生成元集合 $\{a, b\}$ 构造出一个
\newterm{free monoid 自由幺半群}.

\section{Haskell 中的自由幺半群}

在 Haskell 中, 一个二元集合等价于类型 \code{Bool}, 由这个集合生成的自由幺半群等价于类型
\code{{[}Bool{]}} (\code{Bool} 的列表). (我有意忽略无限列表带来的问题.)

Haskell 里的幺半群由这个类型类定义:

\src{snippet01}
它只是说每个 \code{Monoid} 必须有一个中性元素, 叫作 \code{mempty}, 还有一个叫
\code{mappend} 的二元函数 (乘法). 单位律和结合律无法在 Haskell 中表达, 每次实例化一个幺半群时
都必须由程序员自行验证.

任意类型的列表都构成幺半群, 这一事实由这个 instance 定义描述:

\src{snippet02}
它说空列表 \code{{[}{]}} 是单位元素, 列表拼接 \code{(++)} 是二元运算.

正如我们所见, 类型为 \code{a} 的列表对应于以集合 \code{a} 为生成元的自由幺半群.
自然数集合配上乘法不是自由幺半群, 因为我们把大量乘积等同到了一起. 比如对比一下:

\src{snippet03}
这很容易, 但问题是: 我们能不能在范畴论中完成这个自由构造 --- 在那里我们是不允许往物件内部看的?
我们会用我们的老伙计: 普遍构造.

第二个有趣的问题是: 任何一个幺半群是否都能从某个自由幺半群出发、通过比定律所要求的最少数量更多的
元素等同而得到? 我会向你展示, 这一点直接从普遍构造推出.

\section{自由幺半群的普遍构造}

如果你回想一下我们此前与普遍构造打交道的经历, 你可能会注意到, 它与其说是在构造某个东西, 不如说是在
从一个给定模式中挑选最契合的那个物件. 所以, 如果我们想用普遍构造来 ``构造'' 一个自由幺半群, 就必须
先有一大堆幺半群供我们挑选. 我们需要一整个幺半群的范畴可供选择. 但幺半群能构成范畴吗?

我们先把幺半群看作配备了额外结构的集合, 这些结构由单位和乘法定义. 我们挑选那些保持幺半群结构的函数
作为态射. 这种保持结构的函数叫作\newterm{homomorphism 同态}. 幺半群同态必须把两个元素的积映到
这两个元素的像的积:

\src{snippet04}
而且它必须把单位映到单位.

举个例子, 考虑一个从整数列表到整数的同态. 如果把 \code{{[}2{]}} 映到 2, \code{{[}3{]}} 映到
3, 我们就必须把 \code{{[}2, 3{]}} 映到 6, 因为拼接

\src{snippet05}
变成了乘法

\src{snippet06}
现在让我们忘掉单个幺半群的内部结构, 只把它们当作带有相应态射的物件. 你就得到了幺半群的范畴
$\cat{Mon}$.

好吧, 也许在忘掉内部结构之前, 先注意一个重要性质. $\cat{Mon}$ 的每个物件都可以平凡地映到一个
集合 --- 就是它的元素构成的集合. 这个集合叫作\newterm{underlying 底层}集合. 事实上, 我们不仅能
把 $\cat{Mon}$ 的物件映到集合, 还能把 $\cat{Mon}$ 的态射 (同态) 映到函数. 这看上去又有点
平凡, 但马上就会派上用场. 这个从 $\cat{Mon}$ 到 $\Set$ 对物件和态射的映射其实是个函子. 由于这个
函子 ``遗忘'' 了幺半群结构 --- 一旦我们身处一个朴素的集合之中, 就不再区分单位元素、也不关心乘法了
 --- 它叫作\newterm{forgetful functor 遗忘函子}. 遗忘函子在范畴论中经常出现.

我们现在有两种看待 $\cat{Mon}$ 的不同视角. 我们可以把它当作任何其他带有物件和态射的范畴来对待.
在那个视角里, 我们看不到幺半群的内部结构. 关于 $\cat{Mon}$ 中某个特定物件, 我们所能说的一切只是:
它通过态射与自己以及其它物件相连. 态射的 ``乘法'' 表 --- 组合的规则 --- 来自另一个视角:
幺半群即集合. 转入范畴论并没能让我们彻底失去这个视角 --- 我们仍然可以通过遗忘函子访问到它.

要应用普遍构造, 我们需要定义一个特殊性质, 让我们能够在幺半群的范畴中搜寻, 挑出自由幺半群的最佳候选.
但自由幺半群是由它的生成元定义的. 不同的生成元选择产生不同的自由幺半群 (\code{Bool} 的列表和
\code{Int} 的列表不是一回事). 我们的构造必须从一个生成元集合出发. 于是我们又回到了集合!

这正是遗忘函子登场的时候. 我们可以用它给幺半群拍 X 光片. 我们可以在这团团块的 X 光影像里认出生成元.
方法是这样的:

我们从生成元的集合 $x$ 出发. 它是 $\Set$ 中的一个集合.

我们要匹配的模式由一个幺半群 $m$ --- $\cat{Mon}$ 的一个物件 --- 和 $\Set$
中的一个函数 $p$ 构成:

\src{snippet07}
其中 $U$ 是我们从 $\cat{Mon}$ 到 $\Set$ 的遗忘函子. 这是一个古怪的异构模式
--- 一半在 $\cat{Mon}$ 里, 一半在 $\Set$ 里.

思路是: 函数 $p$ 将在 $m$ 的 X 光影像里认出生成元集合. 函数在识别集合中的点方面可能很糟糕
(它们可能把这些点合并掉), 这没关系. 一切都会由普遍构造来摆平, 它会挑出这个模式的最佳代表.

\begin{figure}[H]
  \centering
  \includegraphics[width=0.4\textwidth]{images/monoid-pattern.jpg}
\end{figure}

\noindent
我们还必须定义候选之间的排序. 假设我们有另一个候选: 一个幺半群 $n$ 和一个函数, 在其 X
光影像里认出生成元:

\src{snippet08}
如果存在一个幺半群的态射 (也就是保持结构的同态):

\src{snippet09}
它在 $U$ 下的像 (记住, $U$ 是函子, 所以它把态射映到函数) 经由 $p$ 分解, 我们就说
$m$ 优于 $n$:

\src{snippet10}
如果你把 $p$ 看作在 $m$ 中挑选生成元, 把 $q$ 看作在 $n$ 中挑选 ``同样的''
生成元, 那么你可以把 $h$ 看作在两个幺半群之间映这些生成元. 记住, $h$ 按定义保持幺半群
结构. 这意味着一个幺半群里两个生成元的积, 会被映到第二个幺半群里相应两个生成元的积, 以此类推.

\begin{figure}[H]
  \centering
  \includegraphics[width=0.4\textwidth]{images/monoid-ranking.jpg}
\end{figure}

\noindent
这个排序可以用来找出最佳候选 --- 自由幺半群. 定义如下:

\begin{quote}
  当且仅当存在\emph{唯一}的态射 $h$, 从 $m$ 到任意其他带有函数 $q$ 的幺半群
  $n$, 满足上述分解性质时, 我们说 $m$ (连同函数 $p$) 是以 $x$ 为生成元的
  \textbf{自由幺半群}.
\end{quote}
顺带一提, 这也回答了我们的第二个问题. 函数 $U h$ 正是那个有能力把 $U m$ 的多个元素
合并为 $U n$ 的单个元素的东西. 这种合并对应于把自由幺半群的某些元素等同起来. 因此,
任何以 $x$ 为生成元的幺半群都可以通过把基于 $x$ 的自由幺半群的某些元素等同起来而得到.
自由幺半群是只做了最最少量等同的那个.

谈到伴随的时候, 我们会回到自由幺半群.

\section{挑战}

\begin{enumerate}
  \tightlist
  \item
        你也许会想 (我最初就是这么想的) 幺半群同态必须保持单位这一要求是多余的. 毕竟, 我们知道对一切
        $a$ 都有

        \begin{snip}{text}
h a * h e = h (a * e) = h a
\end{snip}
        所以 $h e$ 表现得像个右单位 (类比地, 也是个左单位). 问题在于, 对全部
        $a$ 的 $h a$ 可能只覆盖目标幺半群的一个子幺半群. 在 $h$ 的像之外,
        可能存在一个 ``真正的'' 单位. 请证明: 保持乘法的幺半群之间的同构必然自动保持单位.
  \item
        考虑一个从带拼接的整数列表到带乘法的整数的幺半群同态. 空列表 \code{{[}{]}} 的像是
        什么? 假设所有单元素列表都被映到它们所含的整数, 即 \code{{[}3{]}} 被映到 3, 等等.
        \code{{[}1, 2, 3, 4{]}} 的像是什么? 有多少个不同的列表被映到整数 12?
        两个幺半群之间还有别的同态吗?
  \item
        由一元集合生成的自由幺半群是什么? 你能看出它与什么同构吗?
\end{enumerate}
